% ===========================================================================
% Modèle de données — les 25 tables, groupées par domaine.
%
% Relevé le 2026-09-22 dans backend/app/models/ (16 modules). Les cardinalités
% proviennent des contraintes réelles : une contrainte d'unicité sur une clé
% étrangère donne 1..1, son absence donne 0..n.
%
% Disposition en BANDES HORIZONTALES, intitulé de domaine dans une gouttière
% à gauche. Les variantes en grille plaçaient les titres de groupe au-dessus
% des entités, qu'ils recouvraient.
%
% Deux faits volontairement visibles parce qu'ils surprennent à la lecture du
% code :
%   - `employees` ne référence PAS `users` : un employé enregistré n'est pas
%     nécessairement titulaire d'un compte ;
%   - `negotiations.candidate_id` est une CHAÎNE, non une clé étrangère : ce
%     lien n'est pas contraint par la base.
% ===========================================================================
\begin{tikzpicture}[
    node distance=7mm and 8mm,
    pqband/.style={
      font=\scriptsize\bfseries, text=piqbitInk, align=right,
      anchor=east, text width=30mm
    }
  ]

  % ═══ Bande 1 — Identité et accès ════════════════════════════════════════
  \node[pqentity, text width=32mm] (users)
    {\textbf{users}\nodepart{two}\texttt{id} · email · rôle · 2FA · état};
  \node[pqentity, right=of users, text width=32mm] (audit)
    {\textbf{audit\_logs}\nodepart{two}acteur · action · cible};
  \node[pqband] at ([xshift=-10mm]users.west) {Identité\\et accès};

  % ═══ Bande 2 — Recrutement ══════════════════════════════════════════════
  \node[pqentity, below=of users, text width=32mm] (cand)
    {\textbf{candidates}\nodepart{two}dossier · CV de référence};
  \node[pqentity, right=of cand, text width=32mm] (jobs)
    {\textbf{job\_offers}\nodepart{two}intitulé · compétences · statut};
  \node[pqentity, right=of jobs, text width=32mm] (apps)
    {\textbf{applications}\nodepart{two}statut · CV soumis · date};
  \node[pqentity, right=of apps, text width=32mm] (saved)
    {\textbf{saved\_jobs}\nodepart{two}favoris du candidat};
  \node[pqband] at ([xshift=-10mm]cand.west) {Recrutement};

  % ═══ Bande 3 — Analyse et notation ══════════════════════════════════════
  \node[pqentity, below=of cand, text width=32mm] (cvan)
    {\textbf{cv\_analyses}\nodepart{two}compétences · expérience};
  \node[pqentity, right=of cvan, text width=32mm] (score)
    {\textbf{candidate\_scores}\nodepart{two}score total · visibilité};
  \node[pqentity, right=of score, text width=32mm] (emb)
    {\textbf{embeddings}\nodepart{two}vecteur 384 dimensions};
  \node[pqentity, right=of emb, text width=32mm] (chunks)
    {\textbf{cv\_chunks}\nodepart{two}segments indexés};
  \node[pqband] at ([xshift=-10mm]cvan.west) {Analyse\\et notation};

  % ═══ Bande 4 — Évaluation ═══════════════════════════════════════════════
  \node[pqentity, below=of cvan, text width=32mm] (inv)
    {\textbf{interview\_invitations}\nodepart{two}jeton public · recruteur};
  \node[pqentity, right=of inv, text width=32mm] (slots)
    {\textbf{interview\_slots}\nodepart{two}début · fin · retenu};
  \node[pqentity, right=of slots, text width=32mm] (tsess)
    {\textbf{test\_sessions}\nodepart{two}jeton · durée · score};
  \node[pqentity, right=of tsess, text width=32mm] (tq)
    {\textbf{test\_questions}\nodepart{two}énoncé · barème};
  \node[pqentity, right=of tq, text width=32mm] (ta)
    {\textbf{test\_answers}\nodepart{two}réponse · correction};
  \node[pqband] at ([xshift=-10mm]inv.west) {Évaluation};

  % ═══ Bande 5 — Embauche ═════════════════════════════════════════════════
  \node[pqentity, below=of inv, text width=32mm] (nego)
    {\textbf{negotiations}\nodepart{two}offre · demande · issue};
  \node[pqentity, right=of nego, text width=32mm] (nround)
    {\textbf{negotiation\_rounds}\nodepart{two}tour · montant · message};
  \node[pqentity, right=of nround, text width=32mm] (contr)
    {\textbf{contracts}\nodepart{two}empreinte · signature};
  \node[pqentity, right=of contr, text width=32mm] (empl)
    {\textbf{employees}\nodepart{two}poste · département · salaire};
  \node[pqband] at ([xshift=-10mm]nego.west) {Embauche};

  % ═══ Bande 6 — RH interne et diffusion ══════════════════════════════════
  \node[pqentity, below=of nego, text width=32mm] (budg)
    {\textbf{department\_budgets}\nodepart{two}enveloppe · exercice};
  \node[pqentity, right=of budg, text width=32mm] (bexp)
    {\textbf{budget\_expenses}\nodepart{two}dépense · imputation};
  \node[pqentity, right=of bexp, text width=32mm] (career)
    {\textbf{career\_plans}\nodepart{two}objectifs · jalons};
  \node[pqentity, right=of career, text width=32mm] (notif)
    {\textbf{notifications}\nodepart{two}type · lu · lien};
  \node[pqband] at ([xshift=-10mm]budg.west) {RH interne\\et diffusion};

  % ═══ Relations principales ══════════════════════════════════════════════
  \draw[pqrel] (users.east)  -- node[pqcard, above] {1..1} (audit.west);
  \draw[pqrel] (users.south) -- node[pqcard, right] {1..1} (cand.north);
  \draw[pqrel] (cand.east)   -- node[pqcard, above] {0..n} (jobs.west);
  \draw[pqrel] (jobs.east)   -- node[pqcard, above] {0..n} (apps.west);
  \draw[pqrel] (apps.east)   -- node[pqcard, above] {0..n} (saved.west);
  \draw[pqrel] (cand.south)  -- node[pqcard, right] {1..1} (cvan.north);
  \draw[pqrel] (cvan.east)   -- node[pqcard, above] {1..1} (score.west);
  \draw[pqrel] (score.east)  -- node[pqcard, above] {0..n} (emb.west);
  \draw[pqrel] (emb.east)    -- node[pqcard, above] {1..n} (chunks.west);
  \draw[pqrel] (cvan.south)  -- node[pqcard, right] {0..n} (inv.north);
  \draw[pqrel] (inv.east)    -- node[pqcard, above] {1..n} (slots.west);
  \draw[pqrel] (slots.east)  -- node[pqcard, above] {0..n} (tsess.west);
  \draw[pqrel] (tsess.east)  -- node[pqcard, above] {0..n} (tq.west);
  \draw[pqrel] (tq.east)     -- node[pqcard, above] {0..n} (ta.west);
  \draw[pqrel] (inv.south)   -- node[pqcard, right] {0..n} (nego.north);
  \draw[pqrel] (nego.east)   -- node[pqcard, above] {1..n} (nround.west);
  \draw[pqrel] (nround.east) -- node[pqcard, above] {0..1} (contr.west);
  \draw[pqrel] (contr.east)  -- node[pqcard, above] {0..1} (empl.west);
  \draw[pqrel] (nego.south)  -- node[pqcard, right] {0..n} (budg.north);
  \draw[pqrel] (budg.east)   -- node[pqcard, above] {1..n} (bexp.west);
  \draw[pqrel] (bexp.east)   -- node[pqcard, above] {0..n} (career.west);
  \draw[pqrel] (career.east) -- node[pqcard, above] {0..n} (notif.west);

  % ═══ Note sur les liens non contraints ══════════════════════════════════
  \node[pqwarn, below=9mm of bexp.south, anchor=north, text width=150mm]
    {\textbf{Deux liens ne sont pas contraints par la base.}
     \texttt{employees} ne porte aucune clé étrangère vers \texttt{users} :
     un employé enregistré n'a pas nécessairement de compte.
     \texttt{negotiations.candidate\_id} est une \emph{chaîne de caractères},
     non une clé étrangère — l'intégrité de ce lien n'est donc pas garantie
     par le moteur. À énoncer comme une dette assumée, non comme un oubli.};

\end{tikzpicture}
